\documentclass[9pt,t]{beamer}
% =====================================================================
% finance-beamer.tex — shared Beamer style for the LLM-in-Finance course
% =====================================================================
% Reusable slide style. Each deck \input's this immediately after
% \documentclass{beamer}. It provides:
%   * the Madrid theme + book colour palette (matches the printed book)
%   * two book-style framing blocks:
%       \begin{bigpicture} ... \end{bigpicture}   ("the bigger picture")
%       \begin{underhood}   ... \end{underhood}    ("under the hood")
%     These mirror the book's `context` and `deepdive` tcolorboxes so the
%     same intuition-vs-mechanism scaffold the reader meets in the book
%     reappears on the slides.
%   * a Python listings style for code frames
%   * appendix conventions: \appendixstart drops a divider and switches
%     to non-numbered backup frames so "extra / less central" material
%     lives after the main talk without inflating the slide count.
%
% Usage:
%   \documentclass{beamer}
%   % =====================================================================
% finance-beamer.tex — shared Beamer style for the LLM-in-Finance course
% =====================================================================
% Reusable slide style. Each deck \input's this immediately after
% \documentclass{beamer}. It provides:
%   * the Madrid theme + book colour palette (matches the printed book)
%   * two book-style framing blocks:
%       \begin{bigpicture} ... \end{bigpicture}   ("the bigger picture")
%       \begin{underhood}   ... \end{underhood}    ("under the hood")
%     These mirror the book's `context` and `deepdive` tcolorboxes so the
%     same intuition-vs-mechanism scaffold the reader meets in the book
%     reappears on the slides.
%   * a Python listings style for code frames
%   * appendix conventions: \appendixstart drops a divider and switches
%     to non-numbered backup frames so "extra / less central" material
%     lives after the main talk without inflating the slide count.
%
% Usage:
%   \documentclass{beamer}
%   % =====================================================================
% finance-beamer.tex — shared Beamer style for the LLM-in-Finance course
% =====================================================================
% Reusable slide style. Each deck \input's this immediately after
% \documentclass{beamer}. It provides:
%   * the Madrid theme + book colour palette (matches the printed book)
%   * two book-style framing blocks:
%       \begin{bigpicture} ... \end{bigpicture}   ("the bigger picture")
%       \begin{underhood}   ... \end{underhood}    ("under the hood")
%     These mirror the book's `context` and `deepdive` tcolorboxes so the
%     same intuition-vs-mechanism scaffold the reader meets in the book
%     reappears on the slides.
%   * a Python listings style for code frames
%   * appendix conventions: \appendixstart drops a divider and switches
%     to non-numbered backup frames so "extra / less central" material
%     lives after the main talk without inflating the slide count.
%
% Usage:
%   \documentclass{beamer}
%   \input{../_style/finance-beamer.tex}
%   \title{...}\subtitle{...}\author{...}\institute{...}\date{\today}
%   \begin{document} ... \end{document}
% =====================================================================

\usetheme{Madrid}
\usecolortheme{whale}
\setbeamertemplate{navigation symbols}{}
\setbeamertemplate{caption}[numbered]
\setbeamercovered{transparent}

% --- Density controls: keep dense, equation-heavy frames on one slide ---
% Slightly smaller body text + tighter lists/blocks. Frame titles and the
% Madrid chrome keep their normal size, so the deck still reads as Madrid,
% just with more vertical room for content.
\setbeamerfont{normal text}{size=\small}
\setbeamertemplate{itemize/enumerate body begin}{\small}
\setbeamertemplate{itemize/enumerate subbody begin}{\footnotesize}
% Tighter block spacing.
\setbeamertemplate{block begin}{%
  \par\vskip2pt%
  \begin{beamercolorbox}[colsep*=.75ex,leftskip=1ex,rightskip=1ex]{block title}%
    \usebeamerfont*{block title}\insertblocktitle%
  \end{beamercolorbox}%
  {\parskip0pt\vskip-1pt}%
  \begin{beamercolorbox}[colsep*=.75ex,vmode,leftskip=1ex,rightskip=1ex]{block body}%
  \vskip1pt}
\setbeamertemplate{block end}{\end{beamercolorbox}\vskip2pt}

\usepackage{amsmath, amssymb}
\usepackage{booktabs}
\usepackage{xcolor}
\usepackage{listings}
\usepackage[skins, breakable]{tcolorbox}

% --- Book colour palette (kept in sync with book/preamble.tex) ---
\definecolor{linkblue}{RGB}{14, 75, 140}
\definecolor{deepdivebg}{RGB}{235, 244, 255}
\definecolor{narrativebg}{RGB}{255, 252, 240}
\definecolor{narrativerule}{RGB}{180, 130, 40}
\definecolor{steelblue}{RGB}{70, 130, 180}
\definecolor{forestgreen}{RGB}{34, 139, 34}
\definecolor{crimson}{RGB}{220, 20, 60}
\definecolor{darkorange}{RGB}{255, 140, 0}
\definecolor{codebg}{RGB}{248, 248, 248}
\definecolor{coderule}{RGB}{210, 210, 210}
\definecolor{keywordblue}{RGB}{0, 100, 180}
\definecolor{commentgray}{RGB}{120, 120, 120}
\definecolor{stringgreen}{RGB}{30, 130, 30}

% Tie the Madrid structural colour to the book's link blue.
\setbeamercolor{structure}{fg=linkblue}
\setbeamercolor{title}{fg=white}
\setbeamercolor{frametitle}{fg=white}

% --- "the bigger picture" : narrative / intuition framing -----------
\newtcolorbox{bigpicture}{%
  enhanced, breakable, boxsep=2pt,
  colback  = narrativebg,
  colframe = narrativerule,
  boxrule  = 0pt, toprule = 1.4pt, bottomrule = 0.4pt,
  leftrule = 0pt, rightrule = 0pt, arc = 0pt,
  left = 6pt, right = 6pt, top = 4pt, bottom = 5pt,
  fonttitle = \footnotesize\scshape\color{narrativerule},
  title = {the bigger picture}, attach title to upper = {\par\smallskip},
}

% --- "under the hood" : the mechanism / technical detail ------------
\newtcolorbox{underhood}{%
  enhanced, breakable, boxsep=2pt,
  colback  = deepdivebg,
  colframe = linkblue,
  boxrule  = 0pt, toprule = 1.4pt, bottomrule = 0.4pt,
  leftrule = 0pt, rightrule = 0pt, arc = 0pt,
  left = 6pt, right = 6pt, top = 4pt, bottom = 5pt,
  fonttitle = \footnotesize\scshape\color{linkblue},
  title = {under the hood}, attach title to upper = {\par\smallskip},
}

% --- Python code style ---------------------------------------------
\lstset{
  basicstyle    = \ttfamily\footnotesize,
  keywordstyle  = \color{keywordblue}\bfseries,
  commentstyle  = \color{commentgray}\itshape,
  stringstyle   = \color{stringgreen},
  showstringspaces = false,
  language      = Python,
  breaklines    = true,
  backgroundcolor = \color{codebg},
  frame         = single,
  rulecolor     = \color{coderule},
  numbers       = none,
  aboveskip     = 4pt, belowskip = 4pt,
}

% --- Appendix conventions ------------------------------------------
% \appendixstart{Title} : begins the "extra material" appendix. Frames
% after it are not counted toward the main deck's frame total, keeping
% the core talk lean while preserving the deeper / optional content.
\newcommand{\appendixstart}[1]{%
  \appendix
  \setbeamertemplate{frametitle continuation}{}%
  \begin{frame}[noframenumbering]
    \begin{center}
      {\Large\scshape\color{linkblue} Appendix}\\[4pt]
      {\large #1}\\[10pt]
      {\footnotesize Extra / optional material — beyond the core lecture.}
    \end{center}
  \end{frame}
}
% Use \begin{frame}[noframenumbering]{...} for each appendix frame.

%   \title{...}\subtitle{...}\author{...}\institute{...}\date{\today}
%   \begin{document} ... \end{document}
% =====================================================================

\usetheme{Madrid}
\usecolortheme{whale}
\setbeamertemplate{navigation symbols}{}
\setbeamertemplate{caption}[numbered]
\setbeamercovered{transparent}

% --- Density controls: keep dense, equation-heavy frames on one slide ---
% Slightly smaller body text + tighter lists/blocks. Frame titles and the
% Madrid chrome keep their normal size, so the deck still reads as Madrid,
% just with more vertical room for content.
\setbeamerfont{normal text}{size=\small}
\setbeamertemplate{itemize/enumerate body begin}{\small}
\setbeamertemplate{itemize/enumerate subbody begin}{\footnotesize}
% Tighter block spacing.
\setbeamertemplate{block begin}{%
  \par\vskip2pt%
  \begin{beamercolorbox}[colsep*=.75ex,leftskip=1ex,rightskip=1ex]{block title}%
    \usebeamerfont*{block title}\insertblocktitle%
  \end{beamercolorbox}%
  {\parskip0pt\vskip-1pt}%
  \begin{beamercolorbox}[colsep*=.75ex,vmode,leftskip=1ex,rightskip=1ex]{block body}%
  \vskip1pt}
\setbeamertemplate{block end}{\end{beamercolorbox}\vskip2pt}

\usepackage{amsmath, amssymb}
\usepackage{booktabs}
\usepackage{xcolor}
\usepackage{listings}
\usepackage[skins, breakable]{tcolorbox}

% --- Book colour palette (kept in sync with book/preamble.tex) ---
\definecolor{linkblue}{RGB}{14, 75, 140}
\definecolor{deepdivebg}{RGB}{235, 244, 255}
\definecolor{narrativebg}{RGB}{255, 252, 240}
\definecolor{narrativerule}{RGB}{180, 130, 40}
\definecolor{steelblue}{RGB}{70, 130, 180}
\definecolor{forestgreen}{RGB}{34, 139, 34}
\definecolor{crimson}{RGB}{220, 20, 60}
\definecolor{darkorange}{RGB}{255, 140, 0}
\definecolor{codebg}{RGB}{248, 248, 248}
\definecolor{coderule}{RGB}{210, 210, 210}
\definecolor{keywordblue}{RGB}{0, 100, 180}
\definecolor{commentgray}{RGB}{120, 120, 120}
\definecolor{stringgreen}{RGB}{30, 130, 30}

% Tie the Madrid structural colour to the book's link blue.
\setbeamercolor{structure}{fg=linkblue}
\setbeamercolor{title}{fg=white}
\setbeamercolor{frametitle}{fg=white}

% --- "the bigger picture" : narrative / intuition framing -----------
\newtcolorbox{bigpicture}{%
  enhanced, breakable, boxsep=2pt,
  colback  = narrativebg,
  colframe = narrativerule,
  boxrule  = 0pt, toprule = 1.4pt, bottomrule = 0.4pt,
  leftrule = 0pt, rightrule = 0pt, arc = 0pt,
  left = 6pt, right = 6pt, top = 4pt, bottom = 5pt,
  fonttitle = \footnotesize\scshape\color{narrativerule},
  title = {the bigger picture}, attach title to upper = {\par\smallskip},
}

% --- "under the hood" : the mechanism / technical detail ------------
\newtcolorbox{underhood}{%
  enhanced, breakable, boxsep=2pt,
  colback  = deepdivebg,
  colframe = linkblue,
  boxrule  = 0pt, toprule = 1.4pt, bottomrule = 0.4pt,
  leftrule = 0pt, rightrule = 0pt, arc = 0pt,
  left = 6pt, right = 6pt, top = 4pt, bottom = 5pt,
  fonttitle = \footnotesize\scshape\color{linkblue},
  title = {under the hood}, attach title to upper = {\par\smallskip},
}

% --- Python code style ---------------------------------------------
\lstset{
  basicstyle    = \ttfamily\footnotesize,
  keywordstyle  = \color{keywordblue}\bfseries,
  commentstyle  = \color{commentgray}\itshape,
  stringstyle   = \color{stringgreen},
  showstringspaces = false,
  language      = Python,
  breaklines    = true,
  backgroundcolor = \color{codebg},
  frame         = single,
  rulecolor     = \color{coderule},
  numbers       = none,
  aboveskip     = 4pt, belowskip = 4pt,
}

% --- Appendix conventions ------------------------------------------
% \appendixstart{Title} : begins the "extra material" appendix. Frames
% after it are not counted toward the main deck's frame total, keeping
% the core talk lean while preserving the deeper / optional content.
\newcommand{\appendixstart}[1]{%
  \appendix
  \setbeamertemplate{frametitle continuation}{}%
  \begin{frame}[noframenumbering]
    \begin{center}
      {\Large\scshape\color{linkblue} Appendix}\\[4pt]
      {\large #1}\\[10pt]
      {\footnotesize Extra / optional material — beyond the core lecture.}
    \end{center}
  \end{frame}
}
% Use \begin{frame}[noframenumbering]{...} for each appendix frame.

%   \title{...}\subtitle{...}\author{...}\institute{...}\date{\today}
%   \begin{document} ... \end{document}
% =====================================================================

\usetheme{Madrid}
\usecolortheme{whale}
\setbeamertemplate{navigation symbols}{}
\setbeamertemplate{caption}[numbered]
\setbeamercovered{transparent}

% --- Density controls: keep dense, equation-heavy frames on one slide ---
% Slightly smaller body text + tighter lists/blocks. Frame titles and the
% Madrid chrome keep their normal size, so the deck still reads as Madrid,
% just with more vertical room for content.
\setbeamerfont{normal text}{size=\small}
\setbeamertemplate{itemize/enumerate body begin}{\small}
\setbeamertemplate{itemize/enumerate subbody begin}{\footnotesize}
% Tighter block spacing.
\setbeamertemplate{block begin}{%
  \par\vskip2pt%
  \begin{beamercolorbox}[colsep*=.75ex,leftskip=1ex,rightskip=1ex]{block title}%
    \usebeamerfont*{block title}\insertblocktitle%
  \end{beamercolorbox}%
  {\parskip0pt\vskip-1pt}%
  \begin{beamercolorbox}[colsep*=.75ex,vmode,leftskip=1ex,rightskip=1ex]{block body}%
  \vskip1pt}
\setbeamertemplate{block end}{\end{beamercolorbox}\vskip2pt}

\usepackage{amsmath, amssymb}
\usepackage{booktabs}
\usepackage{xcolor}
\usepackage{listings}
\usepackage[skins, breakable]{tcolorbox}

% --- Book colour palette (kept in sync with book/preamble.tex) ---
\definecolor{linkblue}{RGB}{14, 75, 140}
\definecolor{deepdivebg}{RGB}{235, 244, 255}
\definecolor{narrativebg}{RGB}{255, 252, 240}
\definecolor{narrativerule}{RGB}{180, 130, 40}
\definecolor{steelblue}{RGB}{70, 130, 180}
\definecolor{forestgreen}{RGB}{34, 139, 34}
\definecolor{crimson}{RGB}{220, 20, 60}
\definecolor{darkorange}{RGB}{255, 140, 0}
\definecolor{codebg}{RGB}{248, 248, 248}
\definecolor{coderule}{RGB}{210, 210, 210}
\definecolor{keywordblue}{RGB}{0, 100, 180}
\definecolor{commentgray}{RGB}{120, 120, 120}
\definecolor{stringgreen}{RGB}{30, 130, 30}

% Tie the Madrid structural colour to the book's link blue.
\setbeamercolor{structure}{fg=linkblue}
\setbeamercolor{title}{fg=white}
\setbeamercolor{frametitle}{fg=white}

% --- "the bigger picture" : narrative / intuition framing -----------
\newtcolorbox{bigpicture}{%
  enhanced, breakable, boxsep=2pt,
  colback  = narrativebg,
  colframe = narrativerule,
  boxrule  = 0pt, toprule = 1.4pt, bottomrule = 0.4pt,
  leftrule = 0pt, rightrule = 0pt, arc = 0pt,
  left = 6pt, right = 6pt, top = 4pt, bottom = 5pt,
  fonttitle = \footnotesize\scshape\color{narrativerule},
  title = {the bigger picture}, attach title to upper = {\par\smallskip},
}

% --- "under the hood" : the mechanism / technical detail ------------
\newtcolorbox{underhood}{%
  enhanced, breakable, boxsep=2pt,
  colback  = deepdivebg,
  colframe = linkblue,
  boxrule  = 0pt, toprule = 1.4pt, bottomrule = 0.4pt,
  leftrule = 0pt, rightrule = 0pt, arc = 0pt,
  left = 6pt, right = 6pt, top = 4pt, bottom = 5pt,
  fonttitle = \footnotesize\scshape\color{linkblue},
  title = {under the hood}, attach title to upper = {\par\smallskip},
}

% --- Python code style ---------------------------------------------
\lstset{
  basicstyle    = \ttfamily\footnotesize,
  keywordstyle  = \color{keywordblue}\bfseries,
  commentstyle  = \color{commentgray}\itshape,
  stringstyle   = \color{stringgreen},
  showstringspaces = false,
  language      = Python,
  breaklines    = true,
  backgroundcolor = \color{codebg},
  frame         = single,
  rulecolor     = \color{coderule},
  numbers       = none,
  aboveskip     = 4pt, belowskip = 4pt,
}

% --- Appendix conventions ------------------------------------------
% \appendixstart{Title} : begins the "extra material" appendix. Frames
% after it are not counted toward the main deck's frame total, keeping
% the core talk lean while preserving the deeper / optional content.
\newcommand{\appendixstart}[1]{%
  \appendix
  \setbeamertemplate{frametitle continuation}{}%
  \begin{frame}[noframenumbering]
    \begin{center}
      {\Large\scshape\color{linkblue} Appendix}\\[4pt]
      {\large #1}\\[10pt]
      {\footnotesize Extra / optional material — beyond the core lecture.}
    \end{center}
  \end{frame}
}
% Use \begin{frame}[noframenumbering]{...} for each appendix frame.


\title{Practical Session 14: Summarizing \& Extracting from SEC Filings}
\subtitle{Hands-On --- Extractive Summarization, ROUGE, and Consistency Checking}
\author{Juan F.\ Imbet}
\institute{Paris Dauphine -- PSL University}
\date{\today}

\begin{document}

\begin{frame}\titlepage\end{frame}

% ============================================================
\begin{frame}{Session Overview}
\textbf{Duration:} 1 hour (50 min problems + 10 min debrief)

\medskip
\begin{block}{Timed agenda}
\begin{enumerate}
  \item \textbf{Recap} (5 min): extractive vs.\ abstractive; ROUGE-$N$; consistency.
  \item \textbf{Problem 1} (12 min): compute ROUGE-1/2/L by hand for two candidate summaries.
  \item \textbf{Problem 2} (13 min): design a faithfulness-checking pipeline; classify
        three summary claims as verified / unverifiable / incorrect.
  \item \textbf{Case Study} (20 min): run the Python 10-K pipeline ---
        section routing, TextRank extractive summary, ROUGE, number consistency.
  \item \textbf{Discussion} (10 min): extractive vs.\ abstractive in production.
\end{enumerate}
\end{block}

\medskip
Code lives in \texttt{code/notebooks/14-financial-text-summarization/}.
\end{frame}

% ============================================================
\begin{frame}{Recap: Formulas You Will Need}
\begin{columns}[T]
\column{0.5\textwidth}
\textbf{ROUGE-$N$ (recall-oriented):}
\[
\frac{\#\{\text{matched } n\text{-grams}\}}{\#\{n\text{-grams in reference}\}}
\]
\textbf{ROUGE-L:} LCS-based; uses the longest common subsequence between candidate
and reference.
\column{0.5\textwidth}
\textbf{Extractive vs.\ abstractive:}
\begin{itemize}
  \item Extractive = verbatim sentences $\to$ faithfulness free.
  \item Abstractive = generated $\to$ fluent but hallucination-prone.
\end{itemize}
\end{columns}

\medskip
\begin{block}{Consistency-check recipe (3 steps)}
(1) Multi-source extraction of each metric $\to$ (2) normalise to common unit/sign
$\to$ (3) cross-reference; flag pairs disagreeing beyond tolerance (e.g.\ 0.1\%).
\end{block}

\medskip
\textbf{Faithfulness label set:} \emph{verified} (matches source) /
\emph{unverifiable} (not in source) / \emph{incorrect} (contradicts source).
\end{frame}

% ============================================================
\begin{frame}{Problem 1: ROUGE by Hand}
\textbf{Reference summary $R$:}\\
\texttt{"revenue rose 13 percent to 56.5 billion"}

\medskip
\textbf{Candidate A:} \texttt{"revenue rose 13 percent"}\\
\textbf{Candidate B:} \texttt{"profit fell 13 percent to 56.5 billion"}

\medskip
\textbf{Tasks (treat each word as a token; lowercase):}
\begin{enumerate}
  \item[\textbf{A.}] Compute \textbf{ROUGE-1} (unigram recall) for A and for B.
  \item[\textbf{B.}] Compute \textbf{ROUGE-2} (bigram recall) for A and for B.
  \item[\textbf{C.}] Compute \textbf{ROUGE-L} (LCS recall) for A and for B.
\end{enumerate}

\medskip
\begin{block}{Discussion question}
Candidate B paraphrases ``revenue rose'' as the \emph{semantically opposite}
``profit fell'' yet shares many tokens. What does its score reveal about ROUGE's
core weakness, and which metric from the lecture would catch the error?
\end{block}
\end{frame}

% ============================================================
\begin{frame}{Problem 2: Designing a Faithfulness Check}
\textbf{Source facts} (from the filing / XBRL):
revenue = \$56.5B (+13\% YoY); net income = \$16.4B; Azure growth = +29\%.

\medskip
\textbf{Candidate abstractive summary:}
\begin{quote}
``Revenue grew 13\% to \$56.5 billion. \textbf{Net income rose to \$17 billion.}
Azure cloud revenue increased \textbf{29\%}, and management \textbf{raised full-year
guidance.}''
\end{quote}

\medskip
\textbf{Tasks:}
\begin{enumerate}
  \item[\textbf{A.}] Label each of the four claims: verified / unverifiable / incorrect.
  \item[\textbf{B.}] For the numerical claims, describe how to ground them against the
        \emph{extracted XBRL} data (which step of the 3-step consistency recipe?).
  \item[\textbf{C.}] Which claim is the most dangerous in a decision context, and why?
\end{enumerate}
\end{frame}

% ============================================================
\begin{frame}{Case Study: An End-to-End 10-K Mini-Pipeline}
\textbf{Goal:} take a 10-K excerpt and run the full chapter pipeline ---
section routing $\to$ extractive summary $\to$ ROUGE vs.\ a reference $\to$
number-consistency check.

\medskip
\textbf{Stack (pure Python, no API key needed):}
\begin{itemize}
  \item \texttt{re} --- section routing \& monetary-figure extraction
  \item \texttt{networkx} --- TextRank graph (sentence similarity $\to$ PageRank)
  \item \texttt{rouge\_score} (or the hand formula) --- evaluation
\end{itemize}

\medskip
\textbf{What you will do:}
\begin{enumerate}
  \item Run the provided code on the sample MD\&A text.
  \item Inspect which sentences TextRank selects.
  \item Run the consistency check; see which prose number disagrees with the table.
\end{enumerate}

\medskip
\begin{alertblock}{Setup}
\texttt{pip install networkx rouge-score} --- no model download or API key required;
everything runs locally in seconds.
\end{alertblock}
\end{frame}

% ============================================================
\begin{frame}[fragile]{Case Study: The Code (Part 1a --- Setup \& Helpers)}
\begin{lstlisting}[basicstyle=\ttfamily\scriptsize]
import re, networkx as nx
from collections import Counter

MDNA = (
 "Revenue increased 13% to $56.5 billion in fiscal 2023. "
 "The growth was driven by strong cloud demand. "
 "Azure revenue grew 29% year over year. "
 "Operating margin expanded to 42%. "
 "Management expects continued momentum into fiscal 2024. "
 "The company faces foreign-exchange headwinds."
)

def sentences(text):
    return [s.strip() for s in re.split(r'(?<=[.!?])\s+', text) if s.strip()]

def bow(s):                       # bag-of-words vector
    return Counter(re.findall(r"[a-z]+", s.lower()))

def sim(a, b):                    # cosine over word counts
    common = set(a) & set(b)
    num = sum(a[w] * b[w] for w in common)
    den = (sum(v*v for v in a.values())**0.5) * \
          (sum(v*v for v in b.values())**0.5)
    return num / den if den else 0.0
\end{lstlisting}
\end{frame}

% ============================================================
\begin{frame}[fragile]{Case Study: The Code (Part 1b --- TextRank)}
\begin{lstlisting}
def textrank(text, k=2):
    sents = sentences(text); vecs = [bow(s) for s in sents]
    G = nx.Graph()
    for i in range(len(sents)):
        for j in range(i+1, len(sents)):
            w = sim(vecs[i], vecs[j])
            if w > 0: G.add_edge(i, j, weight=w)
    scores = nx.pagerank(G)
    top = sorted(scores, key=scores.get, reverse=True)[:k]
    return [sents[i] for i in sorted(top)]

print(textrank(MDNA, k=2))
\end{lstlisting}

\medskip
\textbf{Idea:} sentences vote for each other via cosine similarity; PageRank
surfaces the most central (salient) ones --- selected \emph{verbatim}, so the
extractive summary is faithful by construction.
\end{frame}

% ============================================================
\begin{frame}[fragile]{Case Study: The Code (Part 2 --- Consistency Check)}
\begin{lstlisting}
# Table (XBRL-derived) ground truth, in USD millions
TABLE = {"revenue": 56500.0, "azure_growth_pct": 29.0}

# Numbers stated in the prose narrative
PROSE = {"revenue": 56500.0, "azure_growth_pct": 28.0}  # note: 28 vs 29

def consistency(table, prose, tol=0.001):
    flags = []
    for k in table.keys() & prose.keys():
        t, p = table[k], prose[k]
        rel = abs(t - p) / abs(t) if t else 0.0
        status = "OK" if rel <= tol else "MISMATCH"
        flags.append((k, t, p, f"{rel:.2%}", status))
    return flags

for row in consistency(TABLE, PROSE):
    print(row)
# -> ('revenue', 56500.0, 56500.0, '0.00%', 'OK')
# -> ('azure_growth_pct', 29.0, 28.0, '3.45%', 'MISMATCH')
\end{lstlisting}
\end{frame}

% ============================================================
\begin{frame}{Case Study: Diagnostic Questions}
After running both code parts, work through the following:
\medskip
\begin{enumerate}
  \item Which sentences did TextRank select? Do they cover revenue \emph{and} guidance?
  \item The consistency check flags \texttt{azure\_growth\_pct}. Is this an
        \emph{extraction error} or a \emph{disclosure inconsistency}? How would you tell?
  \item Set \texttt{tol = 0.05}. Does the mismatch still flag? What is the danger of a loose tolerance?
\end{enumerate}
\end{frame}

% ============================================================
\begin{frame}{Discussion: Extractive vs.\ Abstractive in Production}
For each scenario, choose extractive, abstractive, or hybrid --- and justify.

\medskip
\begin{enumerate}
  \item A \textbf{covenant-monitoring} system summarises loan-agreement clauses for a
        credit committee. \emph{Every numerical threshold must be exact.}
  \item An \textbf{equity desk} wants a fluent 150-word daily digest of overnight
        analyst notes for human readers who will not act on it directly.
  \item A \textbf{regulator-facing} 10-K summary that must be auditable: every claim
        traceable to a source sentence.
\end{enumerate}

\medskip
\begin{block}{Frame your answer around}
Faithfulness vs.\ fluency $\cdot$ hallucination cost $\cdot$ auditability /
traceability $\cdot$ whether a human or a downstream system consumes the output
$\cdot$ availability of a structured (XBRL) ground truth for grounding.
\end{block}
\end{frame}

% ============================================================
\begin{frame}{Solution: Problem 1 --- ROUGE by Hand}
Reference $R$ (7 tokens): \texttt{revenue rose 13 percent to 56.5 billion}.

\medskip
\textbf{Candidate A} = \texttt{revenue rose 13 percent}:
\begin{itemize}
  \item ROUGE-1 = $4/7 \approx \mathbf{0.571}$ (revenue, rose, 13, percent matched).
  \item ROUGE-2 = $3/6 = \mathbf{0.500}$ (bigrams ``revenue rose'', ``rose 13'', ``13 percent'').
  \item ROUGE-L = LCS length $4 / 7 \approx \mathbf{0.571}$ (the 4 words appear in order).
\end{itemize}

\medskip
\textbf{Candidate B} = \texttt{profit fell 13 percent to 56.5 billion}:
\begin{itemize}
  \item ROUGE-1 = $5/7 \approx \mathbf{0.714}$ (13, percent, to, 56.5, billion matched).
  \item ROUGE-2 = $4/6 \approx \mathbf{0.667}$ (``13 percent'', ``percent to'', ``to 56.5'', ``56.5 billion'').
  \item ROUGE-L = LCS $= 5/7 \approx \mathbf{0.714}$.
\end{itemize}
\end{frame}

% ============================================================
\begin{frame}{Solution: Problem 1 --- The Lesson}
Recall the scores: Candidate A (\texttt{revenue rose 13 percent}) scored
$\approx 0.571$; Candidate B (\texttt{profit fell 13 percent to 56.5 billion})
scored $\approx 0.714$ on every ROUGE variant.

\medskip
\begin{alertblock}{The lesson}
B \emph{outscores} A on every ROUGE variant despite \textbf{inverting the meaning}
(``revenue rose'' $\to$ ``profit fell''). ROUGE rewards lexical overlap, not
semantics or factual direction --- exactly why you need \textbf{factual-consistency
/ NLI-based} checks (FactCC) for finance.
\end{alertblock}
\end{frame}

% ============================================================
\begin{frame}{Solution: Problem 2 --- Faithfulness Check}
\textbf{A. Claim labels:}
\begin{itemize}
  \item ``Revenue grew 13\% to \$56.5B'' $\to$ \textbf{verified} (matches source).
  \item ``Net income rose to \$17 billion'' $\to$ \textbf{incorrect}
        (source: \$16.4B; rounding to \$17B is a \emph{material} misstatement).
  \item ``Azure increased 29\%'' $\to$ \textbf{verified}.
  \item ``Management raised full-year guidance'' $\to$ \textbf{unverifiable}
        (no guidance fact in the provided source).
\end{itemize}

\medskip
\textbf{B. Grounding numbers against XBRL:} this is step (3),
\emph{cross-reference} --- after multi-source extraction and normalisation,
compare each prose number to the XBRL figure; \$17B vs.\ \$16.4B exceeds a 0.1\%
tolerance ($\sim$3.7\%) $\to$ flag.

\medskip
\textbf{C. Most dangerous:} the \emph{net income} claim --- it is confidently
stated, numerically wrong, and feeds directly into valuation/EPS reasoning. The
``unverifiable'' guidance claim is also risky (potential extrinsic hallucination)
but is detectable as ungrounded rather than contradicted.
\end{frame}

% ============================================================
\begin{frame}{Wrap-Up: Key Takeaways}
\begin{itemize}
  \item \textbf{ROUGE is lexical}: high overlap can coexist with inverted meaning.
        Always pair with semantic (BERTScore) \emph{and} factual (NLI / XBRL-grounded) checks.
  \item \textbf{Extractive summarization} buys faithfulness for free --- TextRank
        selects salient verbatim sentences via sentence-graph PageRank.
  \item \textbf{Consistency checking} is a cheap, high-value quality gate:
        normalise, then cross-reference prose numbers against the XBRL table.
  \item \textbf{Tolerance matters}: too loose a threshold hides material errors;
        too tight flags rounding noise.
  \item \textbf{Choose the paradigm by the stakes}: exact numbers / auditability
        $\to$ extractive; fluent human digests $\to$ abstractive (with checks).
\end{itemize}
\medskip
\begin{block}{Next practical}
Fine-tuning SEC-BERT for FINER-139 NER and benchmarking against a zero-shot LLM
extractor on earnings-call entity recognition.
\end{block}
\end{frame}

% ===========================================================================
% APPENDIX
% ===========================================================================
\appendixstart{Stretch Problems}

\begin{frame}[noframenumbering]{Appendix: Stretch Problem --- Hierarchical S-1}
An S-1 runs 300{,}000 words --- beyond any single context window.

\medskip
\textbf{Task:} sketch a two-stage hierarchical summarizer.
\begin{enumerate}
  \item Stage 1: each of $\sim$30 sections $\to$ 200--400-word summary. What is the
        approximate total intermediate token budget?
  \item Stage 2: concatenate section summaries; second-pass model $\to$ final summary.
  \item Compare cost vs.\ a retrieval-augmented approach that only summarises the
        top-$k$ passages for ``principal risk factors''. When does each win?
\end{enumerate}
\end{frame}

\begin{frame}[noframenumbering]{Appendix: Stretch Problem --- Arithmetic via Code}
The lecture's \emph{ReAct} principle: route arithmetic to a code interpreter.

\medskip
\textbf{Task:} given extracted \texttt{revenue\_2023 = 56500} and
\texttt{revenue\_2022 = 50000} (USD M), an LLM is asked for YoY growth.
\begin{enumerate}
  \item Why might the LLM answer ``$\sim$11\%'' or ``15\%'' when computing inline?
  \item Write the one-line Python the agent should emit instead:
        \texttt{(56500 - 50000) / 50000 * 100}.
  \item State the verified result and explain how the
        ``decide-what-to-compute'' vs.\ ``actually-compute'' separation eliminates
        the numerical error class.
\end{enumerate}
\end{frame}

\end{document}
